% TOP_PROBLEM: {"id": 3047, "title": "Dividing up the unrestricted partitions", "category_id": 2, "category": {"id": 2, "name": "combinatorics", "display_name": "Combinatorics", "description": "Counting problems, graph theory, discrete structures.", "slug": "combinatorics", "order_index": 2, "created_at": "2026-07-31T15:26:25.671Z"}, "status": "open", "problem_number": "OPG-37222", "set_id": 12}
% FIRST_POSTED: 2026-10-01
% ENTRY_KIND: statement_only
% UnsolvedMath ID 3047; source catalogue code OPG-37222
% !TeX program = lualatex
% Definition and sources only; this file is not an LLM solution attempt.
\documentclass[11pt,a4paper]{article}
\usepackage[margin=25mm]{geometry}
\usepackage{fontspec}
\setmainfont{lmroman10-regular.otf}[BoldFont=lmroman10-bold.otf,
 ItalicFont=lmroman10-italic.otf,BoldItalicFont=lmroman10-bolditalic.otf]
\usepackage{amsmath,amssymb,mathtools}
\usepackage{unicode-math}
\setmathfont{latinmodern-math.otf}
\AtBeginDocument{\let\setminus\smallsetminus}
\usepackage{xurl,hyperref}
\hypersetup{colorlinks=true,urlcolor=blue}
\setcounter{secnumdepth}{0}
\setlength{\parindent}{0pt}
\setlength{\parskip}{5pt}
\setlength{\emergencystretch}{3em}
\begin{document}
\section*{Dividing up the unrestricted partitions}\label{3047}
{\small UnsolvedMath ID 3047 · OPG-37222 · Combinatorics · OpenGarden}
\subsection{Definitions and mathematical statement}
Work in the ring $\mathbb Z[[x]]$ of formal power series. For each pair
of positive integers $j,m$, choose a sign $\epsilon_{j,m}\in\{-1,1\}$,
and set $\epsilon_{j,0}=1$. Form the product
\[
 F(x)=\prod_{j=1}^{\infty}\left(\sum_{m=0}^{\infty}
                      \epsilon_{j,m}x^{jm}\right)
      =\sum_{N=0}^{\infty}c_Nx^N.
\]
Each coefficient is well defined: only factors $j\le N$ and terms with
$jm\le N$ can contribute to $c_N$. No analytic convergence condition is
involved, and signs for different pairs $(j,m)$ may be chosen independently.

The question is whether one infinite choice of these signs can make
$c_N\in\{-1,0,1\}$ for every $N\ge0$. More explicitly, if a partition
of $N$ has multiplicities $(m_j)_{j\ge1}$, give it the sign
$\prod_j\epsilon_{j,m_j}$. Then
\[
 c_N=\sum_{\substack{m_j\ge0\text{ finitely supported}\\
                       \sum_jjm_j=N}}\prod_j\epsilon_{j,m_j}.
\]
Thus the same rule for a part size and its multiplicity must balance the
positive and negative partitions at every integer simultaneously. The
coefficient $c_0$ is automatically $1$. Modulo $2$, all signs are $1$, so
$c_N$ has the same parity as the ordinary partition number $p(N)$.
The signs cannot be assigned independently to individual partitions:
their multiplicative factorization through $(j,m)$ is essential. Satisfying
only a finite initial range of coefficients does not itself furnish the
required infinite assignment.
\subsection{Short English statement}
Can signs in the unrestricted-partition product be chosen so every resulting coefficient is zero, one, or minus one?
\subsection{Sources}
\begin{itemize}
\item[S1] ULAM.AI and the UnsolvedMath Contributors, supplied problems.json, record 3047 (OPG-37222), OpenGarden collection. Catalogue identity and source transcription; CC BY 4.0.
\url{https://huggingface.co/datasets/ulamai/UnsolvedMath}.
\item[S2] Open Problem Garden, Dividing up the unrestricted partitions. Original problem and discussion; the mathematical formulation below is expanded from the supplied source record.
\url{http://www.openproblemgarden.org/op/dividing_up_the_unrestricted_partitions}.
\end{itemize}
\end{document}
